\documentclass[10pt,twoside]{article}\usepackage[]{graphicx}\usepackage[dvipsnames,svgnames,table]{xcolor}
% maxwidth is the original width if it is less than linewidth
% otherwise use linewidth (to make sure the graphics do not exceed the margin)
\makeatletter
\def\maxwidth{ %
  \ifdim\Gin@nat@width>\linewidth
    \linewidth
  \else
    \Gin@nat@width
  \fi
}
\makeatother

\definecolor{fgcolor}{rgb}{0.345, 0.345, 0.345}
\newcommand{\hlnum}[1]{\textcolor[rgb]{0.686,0.059,0.569}{#1}}%
\newcommand{\hlsng}[1]{\textcolor[rgb]{0.192,0.494,0.8}{#1}}%
\newcommand{\hlcom}[1]{\textcolor[rgb]{0.678,0.584,0.686}{\textit{#1}}}%
\newcommand{\hlopt}[1]{\textcolor[rgb]{0,0,0}{#1}}%
\newcommand{\hldef}[1]{\textcolor[rgb]{0.345,0.345,0.345}{#1}}%
\newcommand{\hlkwa}[1]{\textcolor[rgb]{0.161,0.373,0.58}{\textbf{#1}}}%
\newcommand{\hlkwb}[1]{\textcolor[rgb]{0.69,0.353,0.396}{#1}}%
\newcommand{\hlkwc}[1]{\textcolor[rgb]{0.333,0.667,0.333}{#1}}%
\newcommand{\hlkwd}[1]{\textcolor[rgb]{0.737,0.353,0.396}{\textbf{#1}}}%
\let\hlipl\hlkwb

\usepackage{framed}
\makeatletter
\newenvironment{kframe}{%
 \def\at@end@of@kframe{}%
 \ifinner\ifhmode%
  \def\at@end@of@kframe{\end{minipage}}%
  \begin{minipage}{\columnwidth}%
 \fi\fi%
 \def\FrameCommand##1{\hskip\@totalleftmargin \hskip-\fboxsep
 \colorbox{shadecolor}{##1}\hskip-\fboxsep
     % There is no \\@totalrightmargin, so:
     \hskip-\linewidth \hskip-\@totalleftmargin \hskip\columnwidth}%
 \MakeFramed {\advance\hsize-\width
   \@totalleftmargin\z@ \linewidth\hsize
   \@setminipage}}%
 {\par\unskip\endMakeFramed%
 \at@end@of@kframe}
\makeatother

\definecolor{shadecolor}{rgb}{.97, .97, .97}
\definecolor{messagecolor}{rgb}{0, 0, 0}
\definecolor{warningcolor}{rgb}{1, 0, 1}
\definecolor{errorcolor}{rgb}{1, 0, 0}
\newenvironment{knitrout}{}{} % an empty environment to be redefined in TeX

\usepackage{alltt}
\usepackage[marginparsep=1em]{geometry}
\geometry{lmargin=1.0in,rmargin=1.0in, bmargin=1.2in,  tmargin=1.2in}
\usepackage[dvipsnames,svgnames,table]{xcolor}
\usepackage{graphicx}
\usepackage{amssymb}
\usepackage{epstopdf}
\usepackage{verbatim}
\usepackage{enumerate}
\usepackage{bm}
\usepackage{amsthm}
\usepackage{float}
\usepackage{amsmath}
\usepackage{fancyhdr}
\usepackage{hyperref}
\usepackage{mathtools}
\usepackage{booktabs}
            
%%%%  SHORTCUT COMMANDS  %%%%
\newcommand{\ds}{\displaystyle}
\newcommand{\Z}{\mathbb{Z}}
\newcommand{\T}{\mathcal{T}}
\newcommand{\arc}{\rightarrow}
\newcommand{\R}{\mathbb{R}}
\newcommand{\RP}{\mathbb{R}(+)}
\newcommand{\Rs}{\mathbb{R}^{**}}
\newcommand{\C}{\mathbb{C}}
\newcommand{\E}{\mathbb{E}}
\newcommand{\B}{\mathcal{B}}
\newcommand{\PX}{\mathcal{P}(X)}
\newcommand*\rot{\rotatebox{90}}
\newcommand*\OK{\ding{51}}
\newcommand{\N}{\mathbb{N}}
\newcommand{\Q}{\mathbb{Q}}
\newcommand{\Answer}{\vspace{2mm}\textbf{\underline{Answer}}\\}
\newcolumntype{L}{>{$}l<{$}}
\newcolumntype{C}{>{$}c<{$}}
\newcommand{\GL}{\text{GL}}
\newcommand{\SL}{\text{SL}}

\newcommand{\stirling}[2]{\genfrac{\{}{\}}{0pt}{}{#1}{#2}}

\pagestyle{fancy}
\fancyhf{}
\renewcommand{\sectionmark}[1]{\markright{\thesection.\ #1}}
\lhead{\fancyplain{}{}} 
\fancyhead[RE,RO]{Stat 5200/6200, Fall 2026}
\fancyfoot[RE,RO]{\thepage}
\IfFileExists{upquote.sty}{\usepackage{upquote}}{}
\begin{document}
\begin{flushright}
\begin{minipage}{.33\textwidth}
\rightline{YOUR NAME}
\rightline{\href{mailto:YOUR EMAIL}{YOUR EMAIL}}
\rightline{\today}
\end{minipage}
\end{flushright}

\begin{center}
{\large{\textbf{Homework 7}}}
\end{center}

\begin{enumerate}
  
    \item (\textbf{Modified Exercise 12.1 from textbook}) Consider a four-factor model with $A$ and $D$ fixed, each with three levels. Factors $B$ and $C$ are random with two levels each. All factors are crossed. There is a total of 72 observations.
  
  \begin{enumerate}
    \item Write an appropriate linear model for this situation. Be careful with subscripts!
    \item How many total parameters need to be estimated?
  \end{enumerate}
  
  \item (\textbf{Modified Exercise 12.2 from textbook}) Consider a four-factor model with $A$ and $D$ fixed, each with three levels. Factors $B$ and $C$ are random with two levels each. B nests in A, C nests in B, and D crosses with the others. There is a total of 72 observations.
  
  \begin{enumerate}
    \item Write an appropriate linear model for this situation. Be careful with subscripts!
    \item How many total parameters need to be estimated?
  \end{enumerate}
  
  \item \textbf{Pig Pens: CRD vs Nesting}.
  The dataset \texttt{pig\_feed\_data.csv} in Canvas contains (synthetic) data from the following experiment:
  
  \begin{quote}
    An agricultural researcher would like to test the impact new feed formulations will have on swine weight.
    They have three formulas they would like to compare, and access to 60 pigs.
    They group together the pigs in pens of 5 pigs each, given 12 total pens.
    They then perform a completely randomized design, assigning each of the three distinct feed formulas to four pens each, and measure the weight change of the pigs after 6 months.
  \end{quote}

The data in Canvas include changes in weight values for every single pig, not just aggregated by pen (as discussed in Module 3).

\begin{enumerate}
  \item \textbf{Fit a fixed-effects only ANOVA model by aggregating the data, as suggested in module 3}.
  That is, we want to avoid pseudo-replication by using only data aggregated by the experimental units (pen).
  The code below can be used to compute the mean-weight for each pen:

\begin{knitrout}
\definecolor{shadecolor}{rgb}{0.969, 0.969, 0.969}\color{fgcolor}\begin{kframe}
\begin{alltt}
\hlcom{# IF NECESSARY: }
\hlcom{# install.packages("dplyr") }
\hlkwd{library}\hldef{(dplyr)}

\hlcom{# Double check your file-paths}
\hldef{pig_data} \hlkwb{<-} \hlkwd{read.csv}\hldef{(}\hlsng{'pig_feed_data.csv'}\hldef{)}

\hlcom{# Get mean-weight for each Pen.}
\hldef{mean_pig_data} \hlkwb{<-} \hldef{pig_data |>}
  \hlkwd{group_by}\hldef{(Pen, Feed) |>}
  \hlkwd{summarize}\hldef{(}\hlkwc{Weight} \hldef{=} \hlkwd{mean}\hldef{(Weight))}
\end{alltt}
\end{kframe}
\end{knitrout}


\item Now we will refit a model to the raw-data (values for each pig) using a nested model structure.

\begin{enumerate}
  \item Write down a model with a nested structure appropriate for these data. Carefully define your subscripts and model assumptions.
  \item Use the \texttt{lme4} and \texttt{lmerTest} packages to fit this model and perform a hypothesis test for the fixed effects. \textbf{You need to load both \texttt{lme4} and \texttt{lmerTest} packages before fitting the model to get the anova() function to work properly}

\begin{knitrout}
\definecolor{shadecolor}{rgb}{0.969, 0.969, 0.969}\color{fgcolor}\begin{kframe}
\begin{alltt}
\hlcom{# IF NECESSARY: }
\hlcom{# install.packages("lme4")}
\hlcom{# install.packages("lmerTest")}

\hlkwd{library}\hldef{(lme4)}
\hlkwd{library}\hldef{(lmerTest)}

\hldef{lmer_mod} \hlkwb{<-} \hlkwd{lmer}\hldef{(...)}

\hlkwd{anova}\hldef{(lmer_mod)}
\end{alltt}
\end{kframe}
\end{knitrout}


\end{enumerate}

\item Compare the F-values and P-values of your two tests.

\item Using either model, perform a post-hoc analysis for the main treatment effect. 

\item (\textbf{Comparing to incorrect approach}) Now fit a fixed-effects model to the raw pig-weight data, treating each measurement as if it were an experimental unit (ignoring the nesting structure).
\textbf{What is the difference in the output, and what could go wrong if we did this in practice?}

\item (\textbf{6000-level only}: What we gain by modeling nesting). One question we might ask: \emph{If we were to repeat the study, what is the optimal allocation of pigs-to-pen, based on budget constraints?} This is actually something that we can compute!

Equation~1 of \texttt{CostSupplement.pdf} (Slide 12) gives the following equation for minimizing the observation variance given cost constraints:

$$
\text{Optimal Number of MUs per EU} = \sqrt{ \frac{\sigma^2_{MU} \cdot \text{(Cost of EU)}}{\sigma^2_{EU} \cdot \text{(Cost of MU)} }}
$$

If a single pig costs $\$75$, and a pen costs $\$500$ to build and maintain, how many pigs-per-pen would your recommend for a follow-up study?


\end{enumerate}

\end{enumerate}


\end{document}
